% Point to the main file
\documentclass[../../Main.tex]{subfiles}

% ========== Start the document ==========
\begin{document}

\section{Section 6.3 - Permuations and Combinations}

Let's start with permutations

\subsection{Permutations}

Suppose you have 8 different books. How many ways can the books on a shelf be arranged if... 

\begin{enumerate}
    \item 
    {
        The shelf fits all 8 books?

        $$\underline{8} \cdot \underline{7} \cdot \underline{6} \cdot \underline{5} \cdot \underline{4} \cdot \underline{3} \cdot \underline{2} \cdot \underline{1} = 8! = \boxed{40320}$$
    }

    \item
    {
        The shelf fits 5 books

        $$\underline{8} \cdot \underline{7} \cdot \underline{6} \cdot \underline{5} \cdot \underline{4} = \frac{8 \cdot 7 \cdot 6 \cdot 5 \cdot 4 \cdot 3 \cdot 2 \cdot 1}{3 \cdot 2 \cdot 1} = \frac{8!}{3!} = \frac{8!}{(8 - 5)!} = \boxed{6720}$$

        Here, the  order matters. We call this a \textbf{permutation}
    }
\end{enumerate}

\begin{definition}[Permutation]
    A \yellowbox{Permutation} is one specific way a group of items can be arranged where order matters. 
    In general, if we have $n$ objects and we want to arrange $k$ of them, then the number of ways is given by

    $$\prescript{}{n}{P}_k = \frac{n!}{(n - k)!}$$
\end{definition}

\begin{example}
    How many ways can the president, vice president, and treasuer be assigned from a group of 4 poeple?

    $$\underbrace{4}_{P} \cdot \underbrace{3}_{VP} \cdot \underbrace{2}_{T} = \boxed{24}$$

    or, using the formula

    $$\prescript{}{4}{P}_3 = \frac{4!}{(4 - 3)!} = \frac{4!}{1!} = \boxed{24}$$
\end{example}

Now let's move on to Combinations

\subsection{Combinations}

\begin{example}
    How many committees of 3 people can be formed from 4 people?

    Here, order \textbf{DOES NOT} matters

    Here, there are 4 different committees

    \begin{tabular}{|c|c|c|c|}
        \hline
        ABC & ABD & ACD & BCD \\
        \hline
    \end{tabular}

    You can arrange the 3 people 6 different ways, given by 

    $$\prescript{}{3}{P}_3 = 3! = 6$$

    \begin{note}
        I would write everything out but I'm too lazy
    \end{note}

    Another way of getting that there are 4 committees is
    
    
    $$\frac{\text{total \# ordered arrangements}}{\text{\# ways of permuting that arrangement}} = \frac{24}{6} = \frac{\frac{4!}{(4 - 3)!}}{3!} = \frac{4!}{3!(4 - 3)!}$$

    This is the number of \textbf{Combinations}
\end{example}

\begin{definition}[Combination]
    A \yellowbox{Combinaton} is just an arrangement where order does not matter.

    The number of combinations of $k$ items taken from $n$ items is

    $${n \choose k} = \prescript{}{n}{C}_k = \frac{n!}{k!(n - k)!}$$

    You would say "n choose k"

    Think of this as "how many ways can you arrange $k$ things in $n$ slots"
\end{definition}

\begin{example}

    \begin{enumerate}
        \item
        {
            How many ways can 2 cards be dealt from a standard deck of cards?

            \begin{center}
                52 cards, want 2 cards
            \end{center}

            \begin{align*}
                {52 \choose 2} &= \frac{52!}{2!{52 - 2}!} \\
                &= \frac{52!}{2! 50!} \\
                &= \frac{52 \cdot 51}{2} \\
                &= \boxed{1326} \\
            \end{align*}
        }

        \item
        {
            A coin is flipped 10 times. How many ways can exactly 4 of the flips be heads?

            \begin{align*}
                {10 \choose 4} &= \frac{10!}{4!6!} \\
                &= \boxed{210} \\
            \end{align*}

            notice how this also works for 6 flips to be tails!
        }

        \item
        {
            Ten distinct points are labeled on a circle. How many quadrilaterals can be drawn using these points?

            $${10 \choose 4} = \frac{10!}{4!6!} = \boxed{210}$$
        }
    \end{enumerate}

    There's a LOT more examples but I'm too lazy to write them all out. Just go back to the notes and write them out later

\end{example}

Let's prove that ${n + 1 \choose k} = {n \choose k - 1} + {n \choose k}$

\begin{poof}
    \begin{align*}
        {n \choose k - 1} + {n \choose k} &= \frac{n!}{(k - 1)!(n - (k - 1))!} + \frac{n!}{k!(n - k)!} \\
        &= 
    \end{align*}

    uhhh do later (it's on the very back of 6.3)
\end{poof}

% ========== End the document ==========
\end{document}  